% Generated draft: requires mathematical and source review.
\documentclass[11pt]{article}
\usepackage[margin=25mm]{geometry}
\usepackage{fontspec}
\setmainfont{Latin Modern Roman}
\usepackage{amsmath,amssymb,mathtools,unicode-math}
\setmathfont{Latin Modern Math}
\usepackage{enumitem,needspace,xurl,hyperref}
\setlength{\parindent}{0pt}
\setlength{\parskip}{5pt}
\setlength{\emergencystretch}{3em}
\newcommand{\N}{\mathbb{N}}
\newcommand{\Z}{\mathbb{Z}}
\newcommand{\Q}{\mathbb{Q}}
\newcommand{\R}{\mathbb{R}}
\newcommand{\C}{\mathbb{C}}
\newcommand{\F}{\mathbb{F}}
\newcommand{\A}{\mathbb{A}}
\newcommand{\PP}{\mathbb{P}}
\newcommand{\E}{\mathbb{E}}
\newcommand{\eps}{\varepsilon}
\newcommand{\dd}{\,\mathrm{d}}
\begin{document}
\section*{UnsolvedMath ID: 700329}
\textbf{Source problem code:} catalogue entry
\subsection{Definitions and mathematical statement}
For a group \(K\), define its set of finite homomorphic images, taken up to isomorphism, by
\[
\mathcal{Q}_{\mathrm{fin}}(K)=\{[K/N]_{\cong}:N\trianglelefteq K,\ [K:N]<\infty\}.
\]
Thus only the isomorphism types of finite quotient groups are recorded, without multiplicities. A group \(K\) is residually finite if, for every \(k\ne 1\), there exist a finite group \(Q\) and a homomorphism \(\varphi:K\to Q\) such that \(\varphi(k)\ne 1\). It is finitely generated if it has a finite generating set.

Let \(F_r\) be the free group of rank \(r\), where \(r\) is a nonnegative integer and \(F_0=\{1\}\). Define its derived series by
\[
F_r^{(0)}=F_r,\qquad F_r^{(j+1)}=[F_r^{(j)},F_r^{(j)}].
\]
For a positive integer \(d\), the free soluble group of rank \(r\) in the variety of groups of derived length at most \(d\) is
\[
S_{r,d}=F_r/F_r^{(d)}.
\]
Here \([L,L]\) denotes the subgroup generated by all commutators \(x^{-1}y^{-1}xy\) with \(x,y\in L\).

\textbf{Questions.} Let \(G\) and \(H\) be arbitrary finitely generated residually finite groups satisfying
\[
\mathcal{Q}_{\mathrm{fin}}(G)=\mathcal{Q}_{\mathrm{fin}}(H).
\]
\begin{enumerate}
\item If at least one of \(G\) and \(H\) is isomorphic to \(F_r\) for some nonnegative integer \(r\), must \(G\cong H\)?
\item If at least one of \(G\) and \(H\) is isomorphic to \(S_{r,d}\) for some nonnegative integer \(r\) and positive integer \(d\), must \(G\cong H\)?
\end{enumerate}
The parenthetical wording ``free (free soluble)'' is interpreted as posing these two separate questions. The source specifies neither ranks nor derived lengths; the definition above uses the standard free object in the variety of soluble groups of bounded derived length. In the second question, the other group is not separately assumed to be soluble.
\subsection{Short English statement}
Must two finitely generated residually finite groups with exactly the same finite quotient groups, up to isomorphism, be isomorphic whenever one is free? The parenthetical variant asks the same question when one is free soluble. The supplied registry labels the problem open; that label is not an independently verified assessment.
\subsection{Sources}
\begin{itemize}
\item Local UnsolvedMath registry, permanent ID \(700329\), exact source code \texttt{catalogue entry}, title ``Kourovka Notebook Problem 5.48.'' The supplied record attributes the problem to V. N. Remeslennikov, proposed in 1976, and labels it open.
\item Supplied source attribution: \emph{The Kourovka Notebook (Unsolved Problems in Group Theory)}, 5th issue (1976), Problem 5.48; also identified with the 21st issue edition (2026), \texttt{arXiv:1401.0300v48}. The supplied collection description names E. I. Khukhro and V. D. Mazurov as editors. These bibliographic details have not been independently verified.
\end{itemize}
\end{document}
